\documentclass[12pt]{article}
\newcommand{\DocShort}{Gráfico simple}
\input{preamble_population.tex}
\begin{document}
\doccover{Método gráfico simple}{2026-05}
\handoutintro{Proyectar la población extendiendo visualmente la tendencia histórica de la curva población--tiempo.}{Cuando se desea una lectura rápida y razonada del patrón general de crecimiento.}

\section{Datos ejemplo}
\begin{center}
\begin{tabular}{lcccccc}
\toprule
Año & 1970 & 1980 & 1990 & 2000 & 2010 & 2020 \\
\midrule
Población & 74,000 & 89,000 & 107,000 & 129,000 & 154,000 & 183,000 \\
\bottomrule
\end{tabular}
\end{center}

\section{Ecuaciones mínimas}
Este método no parte de una ecuación cerrada. La operación consiste en:
\begin{enumerate}
\item graficar la población histórica,
\item trazar una curva suave que represente la tendencia,
\item prolongar la curva hacia el horizonte de diseño,
\item leer los valores proyectados.
\end{enumerate}

\section{Procedimiento}
\begin{enumerate}
\item Ubicar el año en el eje horizontal y la población en el eje vertical.
\item Marcar los puntos históricos.
\item Unirlos con una curva suave sin quiebres artificiales.
\item Extender la curva manteniendo el patrón que muestran los datos.
\item Leer los valores de interés.
\end{enumerate}

\section{Ejemplo resuelto}
\begin{center}
\begin{tikzpicture}
\begin{axis}[
width=0.92\textwidth,
height=7cm,
xlabel={Año},
ylabel={Población},
xmin=1970, xmax=2050,
ymin=50000, ymax=400000,
grid=both,
major grid style={gray!25},
minor grid style={gray!10},
xtick={1970,1980,1990,2000,2010,2020,2030,2040,2050},
ytick={50000,100000,150000,200000,250000,300000,350000,400000},
scaled y ticks=false,
xticklabel style={/pgf/number format/1000 sep={}}
]
\addplot[
color=medblue,
mark=*,
thick,
smooth
] coordinates {
(1970,74000) (1980,89000) (1990,107000) (2000,129000) (2010,154000) (2020,183000)
(2030,216000) (2040,252000) (2050,291000)
};
\addplot[dashed, color=teal] coordinates {(2030,50000) (2030,216000)};
\addplot[dashed, color=teal] coordinates {(2040,50000) (2040,252000)};
\addplot[dashed, color=teal] coordinates {(2050,50000) (2050,291000)};
\node[anchor=south east, fill=white] at (axis cs:2028.8,221500) {216,000};
\node[anchor=south east, fill=white] at (axis cs:2038.8,257500) {252,000};
\node[anchor=south east, fill=white] at (axis cs:2048.8,302500) {291,000};
\end{axis}
\end{tikzpicture}
\end{center}

La lectura guiada del gráfico da:
\[
P_{2030}\approx 216,000,\qquad
P_{2040}\approx 252,000,\qquad
P_{2050}\approx 291,000
\]

\resultbox{La lectura guiada del gráfico sugiere 216,000 habitantes en 2030, 252,000 en 2040 y 291,000 en 2050.}

\section{Teoría breve}
El método gráfico resume la historia completa del crecimiento en una forma visual. Su fortaleza está en captar curvaturas y cambios de pendiente que a veces no quedan bien recogidos por una sola fórmula.

\section{Observaciones de uso}
\begin{itemize}
\item Es un método dependiente del juicio del analista.
\item Conviene usar escalas limpias y una curva suave.
\item Debe contrastarse con métodos numéricos si el proyecto es sensible al caudal de diseño.
\end{itemize}
\end{document}
